\documentclass[a4paper]{article}

\usepackage{graphicx}
\usepackage{amsmath,amssymb}
\usepackage{minted}
\usepackage{multirow}
\usepackage{float}
\usepackage{todonotes}
\usepackage{forest}
\usepackage{xstring}
\usepackage{xcolor}
\usepackage[most]{tcolorbox}
\usepackage{xparse}
\usepackage{natbib}

\usetikzlibrary{graphs,graphdrawing}
\usegdlibrary{layered}

% for external graphviz
\input{/home/donkarlo/Dropbox/repo/nd_ai_project/src/nd_ai/knoweldge_representation/language/formal/computer/programming/tex/latex/code_clips/external_graphviz.tex}

% for tags
\input{/home/donkarlo/Dropbox/repo/nd_ai_project/src/nd_ai/knoweldge_representation/language/formal/computer/programming/tex/latex/parts/control_sequence/kind/semtag/semtag.tex}

% for links
\input{/home/donkarlo/Dropbox/repo/nd_ai_project/src/nd_ai/knoweldge_representation/language/formal/computer/programming/tex/latex/code_clips/seehere.tex}

\title{A Short Introduction to Docker}
\date{}
\author{Mohammad Rahmani}

\begin{document}
    \maketitle
    
    
    \section{What Is Docker?}
        
        Docker is a platform for packaging and running applications inside isolated environments called \emph{containers}. A container includes the application, its dependencies, and the configuration required to run it. This makes the application more consistent across different computers and operating systems.
        
        Unlike a virtual machine, a Docker container does not normally include a complete guest operating system. Containers share the host system's kernel, so they are usually smaller and start faster than virtual machines.
    
    
    \section{Important Concepts}
        
        \begin{itemize}
            \item \textbf{Image:} A read-only template that contains an application and its dependencies.
            \item \textbf{Container:} A running instance of an image.
            \item \textbf{Dockerfile:} A text file containing instructions for building an image.
            \item \textbf{Registry:} A service for storing and distributing Docker images, such as Docker Hub.
        \end{itemize}
    
    
    \section{A Simple Dockerfile}
        
        The following Dockerfile creates a small Python application image:
        
        \begin{minted}{dockerfile}
            FROM python:3.12-slim
            
            WORKDIR /app
            COPY app.py /app/app.py
            
            CMD ["python", "app.py"]
        \end{minted}
        
        The \texttt{FROM} instruction selects a base image. \texttt{WORKDIR} defines the working directory inside the container. \texttt{COPY} adds the application file, and \texttt{CMD} defines the command that runs when the container starts.
    
    
    \section{Basic Commands}
        
        \begin{minted}{bash}
            # Build an image
            docker build -t my-python-app .
            
            # Run a container
            Docker run --name my-container my-python-app
            
            # Show running containers
            Docker ps
            
            # Stop a container
            Docker stop my-container
            
            # Remove a container
            Docker rm my-container
        \end{minted}
        
        Docker is especially useful for development, testing, continuous integration, and deployment because it helps teams run the same application environment in different places.

%        \bibliographystyle{plainnat}
%        \bibliography{/home/donkarlo/Dropbox/repo/nd_shared_research/bibliography/bibliography.bib}
\end{document}
